\section{Von Neumann entropy}

\begin{definition}[Von Neumann entropy]\label{def:von_neumann_entropy}
    \lean{vonNeumannEntropy}
    \leanok
    For a Hermitian matrix $\rho \in \MN{D}$ with eigenvalues
    $\lambda_0, \ldots, \lambda_{D-1}$, the \emph{von Neumann entropy} is
    \[
        S(\rho) = -\sum_{i=0}^{D-1} \lambda_i \log \lambda_i,
    \]
    where $0 \log 0 := 0$.
\end{definition}

\begin{lemma}[Entropy respects equality]\label{lem:entropy_congr}
    \lean{vonNeumannEntropy_congr}
    \leanok
    \uses{def:von_neumann_entropy}
    If two Hermitian matrices are equal, then their von Neumann entropies are
    equal.
\end{lemma}

\begin{proof}\leanok
    Substitute the equality of the matrices in the defining eigenvalue sum.
\end{proof}

\begin{lemma}[Entropy of the zero matrix]\label{lem:entropy_zero}
    \lean{vonNeumannEntropy_zero}
    \leanok
    \uses{def:von_neumann_entropy}
    The zero matrix has zero von Neumann entropy:
    \[
        S(0)=0.
    \]
\end{lemma}

\begin{proof}\leanok
    All eigenvalues of the zero matrix are zero, and $0\log 0=0$.
\end{proof}

\begin{theorem}[Entropy is non-negative]\label{thm:entropy_nonneg}
    \lean{vonNeumannEntropy_nonneg}
    \leanok
    \uses{def:von_neumann_entropy}
    For any density matrix $\rho$, $S(\rho) \ge 0$.
\end{theorem}

\begin{proof}\leanok
    Each eigenvalue $\lambda_i$ of a density matrix satisfies
    $0 \le \lambda_i \le 1$, and $-x \log x \ge 0$ on $[0, 1]$.
    \uses{thm:density_eigenvalues_le_one}
\end{proof}

\begin{lemma}[Eigenvalue sum]\label{lem:density_eigenvalues_sum}
    \lean{densityMatrices_eigenvalues_sum_one}
    \leanok
    The eigenvalues of a density matrix sum to $1$.
\end{lemma}

\begin{proof}\leanok
    Follows from $\tr(\rho) = \sum_i \lambda_i = 1$.
\end{proof}

\begin{theorem}[Eigenvalue bound]\label{thm:density_eigenvalues_le_one}
    \lean{densityMatrices_eigenvalues_le_one}
    \leanok
    \uses{lem:density_eigenvalues_sum}
    Each eigenvalue of a density matrix lies in $[0, 1]$.
\end{theorem}

\begin{proof}\leanok
    Non-negativity comes from positive semidefiniteness. The upper bound
    follows because the eigenvalues are non-negative and sum to $1$.
\end{proof}

\begin{theorem}[Entropy upper bound]\label{thm:entropy_le_log_dim}
    \lean{vonNeumannEntropy_le_log_dim}
    \leanok
    \uses{def:von_neumann_entropy}
    For a density matrix $\rho \in \MN{D}$ with $D \ge 1$,
    \[
        S(\rho) \le \log D.
    \]
\end{theorem}

\begin{proof}\leanok
    By Jensen's inequality applied to the concave function $-x \log x$,
    the entropy is maximized when all eigenvalues equal $1/D$.
    \uses{lem:density_eigenvalues_sum, thm:density_eigenvalues_le_one}
\end{proof}

\begin{theorem}[Entropy bounded by log of rank]\label{thm:entropy_le_log_rank}
    \lean{vonNeumannEntropy_le_log_rank}
    \leanok
    \uses{def:von_neumann_entropy}
    For a density matrix $\rho$ of rank $r$,
    \[
        S(\rho) \le \log r.
    \]
    This refines the dimension bound: only the nonzero eigenvalues contribute to
    the entropy, and there are exactly $r$ of them.
\end{theorem}

\begin{proof}\leanok
    The entropy is the $-x\log x$ sum over all eigenvalues; the $D-r$ zero
    eigenvalues contribute nothing. Jensen's inequality applied to $-x\log x$
    over the $r$ nonzero eigenvalues, with uniform weights $1/r$, gives
    $S(\rho) \le \log r$, the maximum attained when each nonzero eigenvalue
    equals $1/r$. The rank equals the number of nonzero eigenvalues of the
    Hermitian matrix $\rho$.
    \uses{lem:density_eigenvalues_sum, thm:density_eigenvalues_le_one}
\end{proof}

\begin{lemma}[Entropy from characteristic polynomial roots]
    \label{lem:entropy_charpoly_roots}
    \lean{vonNeumannEntropy_eq_charpoly_roots}
    \leanok
    \uses{def:von_neumann_entropy}
    The von Neumann entropy of a Hermitian matrix is the $-x\log x$ sum over the
    real parts of the roots of its characteristic polynomial $\chi_\rho$:
    \[
        S(\rho) = \sum_{\lambda \in \mathrm{roots}(\chi_\rho)}
        \bigl({-}\re(\lambda)\, \log \re(\lambda)\bigr).
    \]
\end{lemma}

\begin{proof}\leanok
    The eigenvalues of a Hermitian matrix are exactly the roots of its
    characteristic polynomial, counted with multiplicity, so the eigenvalue sum
    defining $S(\rho)$ equals the displayed sum over roots.
\end{proof}

\begin{theorem}[Trace-logarithm form of the entropy]
    \label{thm:entropy_eq_neg_trace_mul_log}
    \lean{vonNeumannEntropy_eq_neg_trace_mul_log}
    \leanok
    \uses{def:von_neumann_entropy}
    Let $\rho$ be a Hermitian matrix and let $\log\rho$ be its logarithm defined
    through the functional calculus. Then
    \[
        S(\rho) = -\re\bigl(\tr(\rho\,\log\rho)\bigr).
    \]
\end{theorem}

\begin{proof}\leanok
    \uses{def:von_neumann_entropy}
    A Hermitian matrix and its logarithm are simultaneously diagonalized by a
    unitary $U$ with $\rho = U\diag(\lambda_i)U^*$, so
    \[
        \tr(\rho\,\log\rho)
            = \tr\!\bigl(U\diag(\lambda_i\log\lambda_i)U^*\bigr)
            = \sum_i \lambda_i\log\lambda_i.
    \]
    Its negative is $\sum_i\bigl({-}\lambda_i\log\lambda_i\bigr) = S(\rho)$. Zero
    eigenvalues contribute nothing under the convention $0\log 0 = 0$, so no
    full-support assumption is required.
\end{proof}

\begin{remark}
    The logarithm is the totalized real logarithm, with
    $\log x = \log\lvert x\rvert$ and $\log 0 = 0$; both sides of the identity
    use it on every eigenvalue, so the equality holds for an arbitrary Hermitian
    matrix. It coincides with the physical entropy $-\tr(\rho\log\rho)$
    precisely when $\rho$ is positive semidefinite. In that case (a density
    matrix $\rho$, positive semidefinite with unit trace) the eigenvalues
    $\lambda_i$ are nonnegative, $\rho\,\log\rho$ is Hermitian with real trace,
    and the real-part extraction is superfluous: the identity reduces to the
    standard expression $S(\rho) = -\tr(\rho\log\rho)$.
\end{remark}

\begin{definition}[Quantum relative entropy]\label{def:quantum_relative_entropy}
    \lean{quantumRelativeEntropy}
    \leanok
    For matrices $\rho,\sigma\in \MN{D}$, define the trace-log expression
    \[
        D(\rho\Vert\sigma)
            = \re\tr\bigl(\rho(\log\rho-\log\sigma)\bigr).
    \]
    On the physical domain where $\rho$ is a density matrix and $\sigma$ is
    positive definite, this is the Umegaki relative entropy.
\end{definition}

\begin{lemma}[Trace-log splitting for relative entropy]
    \label{lem:quantum_relative_entropy_trace_split}
    \lean{quantumRelativeEntropy_eq_trace_mul_log_sub}
    \leanok
    \uses{def:quantum_relative_entropy}
    For matrices $\rho,\sigma\in \MN{D}$,
    \[
        D(\rho\Vert\sigma)
            = \re\tr(\rho\log\rho)
              - \re\tr(\rho\log\sigma).
    \]
\end{lemma}

\begin{proof}\leanok
    Expand the matrix product over the difference $\log\rho-\log\sigma$ and use
    linearity of the trace and of the real part.
\end{proof}

\begin{lemma}[Relative entropy with identical arguments]
    \label{lem:quantum_relative_entropy_self}
    \lean{quantumRelativeEntropy_self}
    \leanok
    \uses{def:quantum_relative_entropy}
    For every matrix $\rho\in \MN{D}$,
    \[
        D(\rho\Vert\rho)=0.
    \]
\end{lemma}

\begin{proof}\leanok
    The logarithmic difference $\log\rho-\log\rho$ vanishes.
\end{proof}

\begin{lemma}[Relative entropy with zero first argument]
    \label{lem:quantum_relative_entropy_zero_left}
    \lean{quantumRelativeEntropy_zero_left}
    \leanok
    \uses{def:quantum_relative_entropy}
    For every matrix $\sigma\in \MN{D}$,
    \[
        D(0\Vert\sigma)=0.
    \]
\end{lemma}

\begin{proof}\leanok
    The trace-log expression is multiplied on the left by the zero matrix.
\end{proof}

\begin{theorem}[Relative entropy in entropy form]
    \label{thm:quantum_relative_entropy_entropy_form}
    \lean{quantumRelativeEntropy_eq_neg_entropy_sub_trace_mul_log}
    \leanok
    \uses{def:quantum_relative_entropy, lem:quantum_relative_entropy_trace_split,
        thm:entropy_eq_neg_trace_mul_log}
    If $\rho$ is Hermitian, then
    \[
        D(\rho\Vert\sigma)
            = -S(\rho)-\re\tr(\rho\log\sigma).
    \]
\end{theorem}

\begin{proof}\leanok
    Apply Lemma~\ref{lem:quantum_relative_entropy_trace_split} to write
    \[
        D(\rho\Vert\sigma)
            = \re\tr(\rho\log\rho)
              - \re\tr(\rho\log\sigma).
    \]
    The trace-logarithm identity
    \[
        \re\tr(\rho\log\rho)=-S(\rho)
    \]
    then gives the result.
\end{proof}

\begin{theorem}[Klein's inequality]\label{thm:klein_inequality}
    \lean{quantumRelativeEntropy_nonneg}
    \leanok
    \uses{def:quantum_relative_entropy, lem:quantum_relative_entropy_trace_split}
    Let $\rho,\sigma\in \MN{D}$ be density matrices with $\sigma$ of full rank.
    Then the relative entropy is nonnegative,
    \[
        D(\rho\Vert\sigma)\ge 0.
    \]
\end{theorem}

\begin{proof}\leanok
    Diagonalize $\rho=\sum_i p_i\,\ketbra{e_i}{e_i}$ and
    $\sigma=\sum_j q_j\,\ketbra{f_j}{f_j}$ in their eigenbases, with all
    $q_j>0$ since $\sigma$ has full rank. The overlap numbers
    $P_{ij}=\lvert\braket{e_i}{f_j}\rvert^2$ are nonnegative with row sums and column
    sums equal to $1$, because the two eigenbases are orthonormal. A trace
    computation gives
    \[
        \re\tr(\rho\log\rho)=\sum_i p_i\log p_i,
        \qquad
        \re\tr(\rho\log\sigma)
            =\sum_{i,j} p_i\,P_{ij}\log q_j,
    \]
    so that
    \[
        D(\rho\Vert\sigma)
            =\sum_i p_i\log p_i-\sum_{i,j} p_i\,P_{ij}\log q_j.
    \]
    The row sums let the first term be written over the index pair, and the
    scalar tangent bound $\log x\le x-1$ applied to $q_j/p_i$ gives, after
    weighting by $p_i\,P_{ij}$ and summing, the lower bound
    $\sum_{i,j} P_{ij}(q_j-p_i)$. By the column sums and the trace-one
    normalizations of $p$ and $q$ this telescopes to
    $\sum_j q_j-\sum_i p_i=0$, which is the claim.
\end{proof}

\begin{theorem}[Klein's inequality, support form]\label{thm:klein_inequality_general}
    \lean{quantumRelativeEntropy_nonneg_general}
    \leanok
    \uses{def:quantum_relative_entropy}
    Let $\rho,\sigma\in \MN{D}$ be density matrices satisfying the support
    condition $\ker\sigma\subseteq\ker\rho$, that is, every vector annihilated by
    $\sigma$ is annihilated by $\rho$. Then the relative entropy is nonnegative,
    \[
        D(\rho\Vert\sigma)\ge 0.
    \]
\end{theorem}

\begin{proof}\leanok
    \uses{thm:klein_inequality}
    Regularize $\sigma$ by the trace-one perturbation
    $\sigma_\varepsilon' = (1+\varepsilon D)^{-1}(\sigma+\varepsilon\,\mathbf 1)$,
    which is positive definite for every $\varepsilon>0$, hence of full rank. The
    full-rank Klein inequality (Theorem~\ref{thm:klein_inequality}) gives
    $D(\rho\Vert\sigma_\varepsilon')\ge 0$. The perturbation shares the
    eigenbasis $\{\ket{f_j}\}$ of $\sigma$, so the cross term reduces to a scalar
    sum over the eigenvalues $q_j$ of $\sigma$,
    \[
        \re\tr\bigl(\rho\log\sigma_\varepsilon'\bigr)
          = \sum_j \bra{f_j}\rho\ket{f_j}\,
            \log\!\bigl((1+\varepsilon D)^{-1}(q_j+\varepsilon)\bigr).
    \]
    For $q_j>0$ the scalar logarithm converges to $\log q_j$; for $q_j=0$ the
    eigenvector $\ket{f_j}$ lies in $\ker\sigma\subseteq\ker\rho$, so its diagonal
    weight $\bra{f_j}\rho\ket{f_j}$ vanishes and the summand is identically zero.
    No eigenvalue-continuity input is needed, so $D(\rho\Vert\sigma_\varepsilon')
    \to D(\rho\Vert\sigma)$ as $\varepsilon\to0^+$ and the inequality passes to
    the limit.
\end{proof}

\begin{theorem}[Equality case of Klein's inequality]
    \label{thm:klein_inequality_equality}
    \lean{quantumRelativeEntropy_eq_zero_iff}
    \leanok
    \uses{thm:klein_inequality}
    Let $\rho,\sigma\in \MN{D}$ be density matrices with $\sigma$ of full rank.
    Then the relative entropy vanishes exactly when the states coincide,
    \[
        D(\rho\Vert\sigma) = 0 \iff \rho = \sigma.
    \]
    Together with nonnegativity, this is the order property that makes $D$ a
    divergence.
\end{theorem}

\begin{proof}\leanok
    \uses{lem:quantum_relative_entropy_self}
    If $\rho=\sigma$ then $\log\rho-\log\sigma=0$ and $D(\rho\Vert\sigma)=0$.
    For the converse, diagonalize $\rho=\sum_i p_i\,\ketbra{e_i}{e_i}$ and
    $\sigma=\sum_j q_j\,\ketbra{f_j}{f_j}$, with all $q_j>0$ by full rank, and set
    $P_{ij}=\lvert\braket{e_i}{f_j}\rvert^2$. As in
    Theorem~\ref{thm:klein_inequality},
    \[
        D(\rho\Vert\sigma)
            =\sum_{i,j} P_{ij}\bigl[p_i\log p_i-p_i\log q_j-(p_i-q_j)\bigr],
    \]
    a sum of nonnegative terms: each is bounded below by the tangent inequality
    $\log x\le x-1$ at $x=q_j/p_i$, while the linear remainder telescopes through
    the doubly stochastic row and column sums,
    \[
        \sum_{i,j} P_{ij}(p_i-q_j) = \sum_i p_i - \sum_j q_j = 0.
    \]
    If the total vanishes then every term vanishes. On a row with $p_i=0$ the term
    reads $P_{ij}q_j$, so $P_{ij}=0$; on a row with $p_i>0$ the term forces the
    tangent inequality to be tight,
    \[
        \log\!\frac{q_j}{p_i} = \frac{q_j}{p_i}-1,
    \]
    and strict concavity of the logarithm ($\log x<x-1$ for $x\neq1$) gives
    $q_j=p_i$. Hence $q_j=p_i$ whenever $\braket{e_i}{f_j}\neq0$. Writing
    $W=U_\rho^\dagger U_\sigma$ for the overlap of the eigenvector unitaries, this
    matching says $W$ intertwines the eigenvalue diagonals,
    \[
        W\,\mathrm{diag}(q) = \mathrm{diag}(p)\,W,
    \]
    so conjugating $\sigma$ into the eigenbasis of $\rho$ and using the unitarity
    $WW^\dagger=1$,
    \[
        U_\rho^\dagger\sigma U_\rho
            = W\,\mathrm{diag}(q)\,W^\dagger
            = \mathrm{diag}(p)\,WW^\dagger
            = \mathrm{diag}(p),
    \]
    the spectral diagonal of $\rho$. Therefore $\sigma = U_\rho\,\mathrm{diag}(p)\,
    U_\rho^\dagger = \rho$.
\end{proof}

\begin{theorem}[Joint convexity of the relative entropy]
    \label{thm:relative_entropy_joint_convexity}
    \lean{TNLean.RelativeEntropyConvexity.convexOn_quantumRelativeEntropy}
    \leanok
    \uses{def:quantum_relative_entropy, cor:lieb_concavity_id}
    On pairs of positive definite matrices in $\MN{D}$, the map
    $(\rho,\sigma)\mapsto D(\rho\Vert\sigma)$ is jointly convex.
\end{theorem}

\begin{proof}\leanok
    \uses{cor:lieb_concavity_id}
    For $s\in[0,1)$ consider the approximant
    \[
        g_s(\rho,\sigma)
            =\frac{\tr\rho
                -\re\tr(\rho^s\sigma^{1-s})}{1-s}.
    \]
    The trace $\tr\rho$ is real-affine in the pair, hence convex,
    while $(\rho,\sigma)\mapsto\re\tr(\rho^s
    \sigma^{1-s})$ is jointly concave by the $K=\Id$ case of the Lieb concavity
    theorem (Corollary~\ref{cor:lieb_concavity_id}). Their difference, scaled by
    the nonnegative factor $(1-s)^{-1}$, is therefore jointly convex, so each
    $g_s$ is convex.

    Writing $\rho$ and $\sigma$ in their eigenbases with eigenvalues $p_i,q_j>0$
    and overlap weights $P_{ij}=\lvert\braket{e_i}{f_j}\rvert^2$, the approximant
    becomes the double sum
    $g_s(\rho,\sigma)=\sum_{i,j}(1-s)^{-1}(p_i-p_i^s q_j^{1-s})\,P_{ij}$. As
    $s\to1^-$ each per-pair term converges, via the scalar limit
    $(c^u-1)/u\to\log c$, to $p_i(\log p_i-\log q_j)\,P_{ij}$, whose sum equals
    $D(\rho\Vert\sigma)$. Thus $g_s\to D$ pointwise on positive definite pairs,
    and since the pointwise limit of convex functions is convex, $D$ is jointly
    convex.
\end{proof}

\begin{theorem}[Joint convexity on the support domain]
    \label{thm:relative_entropy_joint_convexity_support}
    \lean{TNLean.RelativeEntropyConvexity.convexOn_quantumRelativeEntropy_support}
    \leanok
    \uses{def:quantum_relative_entropy, thm:relative_entropy_joint_convexity}
    On pairs of density matrices $(\rho,\sigma)$ in $\MN{D}$ satisfying the support
    condition $\ker\sigma\subseteq\ker\rho$, the map
    $(\rho,\sigma)\mapsto D(\rho\Vert\sigma)$ is jointly convex.
\end{theorem}

\begin{proof}\leanok
    \uses{thm:relative_entropy_joint_convexity}
    The domain is convex: for a strict convex combination of two such pairs the
    kernel of $a\sigma_1+b\sigma_2$ is $\ker\sigma_1\cap\ker\sigma_2$, since the
    quadratic forms of the positive semidefinite summands are nonnegative and
    their positively weighted sum vanishes only when each does, and a positive
    semidefinite matrix annihilates exactly the vectors of zero quadratic form;
    the two pointwise support inclusions then give
    $\ker\sigma_1\cap\ker\sigma_2\subseteq\ker\rho_1\cap\ker\rho_2$.

    Regularize both arguments through the affine trace-one perturbation
    $M_\varepsilon = (1+\varepsilon N)^{-1}(M+\varepsilon\,\mathbf 1)$, where $N$
    is the matrix dimension, which is positive definite for every $\varepsilon>0$.
    Because the perturbation is affine, it commutes with the convex combination,
    so the four regularized endpoints and the regularized mixture all lie in the
    positive definite domain and the positive definite joint convexity
    (Theorem~\ref{thm:relative_entropy_joint_convexity}) gives the two-point
    inequality for the regularized pairs. As $\varepsilon\to0^+$ the regularized
    relative entropy converges on each pair of the support domain,
    \[
        D(\rho_\varepsilon\Vert\sigma_\varepsilon) \to D(\rho\Vert\sigma).
    \]
    Indeed the perturbation shares the eigenbasis of its argument, so each
    trace-logarithm term is the diagonal sum
    $\sum_j w_j(\varepsilon)\,\log\bigl((1+\varepsilon N)^{-1}(q_j+\varepsilon)\bigr)$,
    where $q_j$ runs over the eigenvalues of $\sigma$ and $w_j(\varepsilon)$ is the
    corresponding diagonal weight. For a positive eigenvalue the scalar factor has
    the limit
    \[
        \log\bigl((1+\varepsilon N)^{-1}(q_j+\varepsilon)\bigr) \to \log q_j
        \qquad (q_j>0),
    \]
    while at a zero eigenvalue the support condition makes the weight vanish, with
    $w_j(\varepsilon)=(1+\varepsilon N)^{-1}\varepsilon$, so the summand is
    $(1+\varepsilon N)^{-1}\varepsilon\,\log\bigl((1+\varepsilon N)^{-1}\varepsilon\bigr)\to0$
    by $x\log x\to0$ as $x\to0^+$. The two-point inequality therefore passes to the
    limit, giving joint convexity on the support domain.
\end{proof}

\begin{lemma}[Functional calculus under unitary conjugation]
    \label{lem:cfc_conj_unitary}
    \lean{Matrix.cfc_conj_unitary}
    \leanok
    For a Hermitian matrix $A$, a real function $f$, and a unitary $U$, the
    continuous functional calculus satisfies
    \[
        f(U A U^\dagger) = U\,f(A)\,U^\dagger.
    \]
\end{lemma}

\begin{proof}\leanok
    The conjugation $\varphi : x \mapsto U x U^\dagger$ is a star-algebra
    automorphism of the matrix algebra, and the continuous functional calculus
    commutes with such automorphisms, $\varphi(f(A)) = f(\varphi(A))$; since
    $\varphi(A) = U A U^\dagger$, this is the claim.
\end{proof}

\begin{lemma}[Matrix logarithm under unitary conjugation]
    \label{lem:log_conj_unitary}
    \lean{Matrix.log_conj_unitary}
    \leanok
    For a Hermitian matrix $A$ and a unitary $U$,
    \[
        \log(U A U^\dagger) = U\,(\log A)\,U^\dagger.
    \]
\end{lemma}

\begin{proof}\leanok
    \uses{lem:cfc_conj_unitary}
    This is the $f = \log$ case of Lemma~\ref{lem:cfc_conj_unitary}: with
    $f(U A U^\dagger) = U\,f(A)\,U^\dagger$ specialized to the real logarithm,
    \[
        \log(U A U^\dagger) = U\,(\log A)\,U^\dagger.
    \]
\end{proof}

\begin{theorem}[Unitary invariance of the quantum relative entropy]
    \label{thm:relative_entropy_unitary_invariance}
    \lean{quantumRelativeEntropy_conj_unitary}
    \lean{Matrix.quantumRelativeEntropy_weyl_conj_eq}
    \leanok
    \uses{def:quantum_relative_entropy}
    For Hermitian matrices $\rho, \sigma$ and a unitary $U$,
    \[
        D(U\rho U^\dagger \,\Vert\, U\sigma U^\dagger) = D(\rho\Vert\sigma).
    \]
\end{theorem}

\begin{proof}\leanok
    \uses{lem:log_conj_unitary}
    Carry the two logarithms through the conjugation by
    Lemma~\ref{lem:log_conj_unitary}, so that
    \[
        D(U\rho U^\dagger \,\Vert\, U\sigma U^\dagger)
            = \re\tr\!\bigl(
                U\rho(\log\rho-\log\sigma)U^\dagger\bigr)
            = \re\tr\!\bigl(
                \rho(\log\rho-\log\sigma)\bigr)
            = D(\rho\Vert\sigma),
    \]
    where the middle equality is trace cyclicity together with $U^\dagger U = 1$.
\end{proof}

\begin{lemma}[Logarithm of a tensor product of positive definite matrices]
    \label{lem:log_kronecker}
    \lean{Matrix.log_kronecker}
    \leanok
    For positive definite matrices $\rho$ and $\tau$,
    \[
        \log(\rho \otimes \tau)
            = \log\rho \otimes \mathbf 1 + \mathbf 1 \otimes \log\tau.
    \]
\end{lemma}

\begin{proof}\leanok
    The tensor product factors as
    $\rho \otimes \tau = (\rho \otimes \mathbf 1)(\mathbf 1 \otimes \tau)$ into a
    pair of commuting positive definite matrices, so the logarithm of the product
    is the sum of the logarithms of the factors. Each unital embedding
    $A \mapsto A \otimes \mathbf 1$ and $B \mapsto \mathbf 1 \otimes B$ is a
    continuous star-algebra homomorphism, hence commutes with the functional
    calculus, which carries each logarithm onto its factor:
    $\log(\rho \otimes \mathbf 1) = \log\rho \otimes \mathbf 1$ and
    $\log(\mathbf 1 \otimes \tau) = \mathbf 1 \otimes \log\tau$.
\end{proof}

\begin{theorem}[Ancilla additivity of the quantum relative entropy]
    \label{thm:relative_entropy_ancilla_additivity}
    \lean{quantumRelativeEntropy_kronecker}
    \leanok
    \uses{def:quantum_relative_entropy, lem:log_kronecker}
    For positive definite matrices $\rho, \sigma$ and a positive definite matrix
    $\tau$ of unit trace,
    \[
        D(\rho \otimes \tau \,\Vert\, \sigma \otimes \tau) = D(\rho\Vert\sigma).
    \]
\end{theorem}

\begin{proof}\leanok
    \uses{lem:log_kronecker}
    Splitting each tensor logarithm by Lemma~\ref{lem:log_kronecker}, the common
    $\mathbf 1 \otimes \log\tau$ terms cancel in the difference, leaving
    $\log(\rho \otimes \tau) - \log(\sigma \otimes \tau)
        = (\log\rho - \log\sigma) \otimes \mathbf 1$. Hence
    \[
        D(\rho \otimes \tau \,\Vert\, \sigma \otimes \tau)
            = \re\tr\!\bigl(
                (\rho(\log\rho - \log\sigma)) \otimes \tau\bigr)
            = \re\bigl(
                \tr(\rho(\log\rho - \log\sigma))
                \cdot \tr\tau\bigr),
    \]
    where the trace of a tensor product factors as a product of traces. As
    $\tr\tau = 1$, the right-hand side is $D(\rho\Vert\sigma)$.
\end{proof}

\begin{lemma}[Root-of-unity character-sum orthogonality]
    \label{lem:root_of_unity_orthogonality}
    \lean{Matrix.sum_rootOfUnity_pow_eq}
    \leanok
    Let $\zeta$ be a primitive $d$-th root of unity and let $i, j$ range over the
    residues modulo $d$. Then
    \[
        \sum_{b=0}^{d-1} \zeta^{b\,i}\,\overline{\zeta^{b\,j}}
            = \begin{cases} d & i = j, \\ 0 & i \ne j. \end{cases}
    \]
\end{lemma}

\begin{proof}\leanok
    Each summand equals $\xi^{b}$ with
    $\xi = \zeta^{i}\,\overline{\zeta}^{\,j}$, and $\xi^{d} = 1$. When $i = j$ the
    base $\xi$ equals $1$ and the sum is $d$. When $i \ne j$ the base is a root of
    unity different from $1$, so the geometric sum
    $(\xi - 1)\sum_{b} \xi^{b} = \xi^{d} - 1 = 0$ forces the sum to vanish; the
    equivalence $\xi = 1 \iff i = j$ uses that $\overline{\zeta} = \zeta^{-1}$ and
    the injectivity of $b \mapsto \zeta^{b}$ on residues.
\end{proof}

\begin{theorem}[Unitarity of the Weyl operators]
    \label{thm:weyl_operator_unitary}
    \lean{Matrix.weylShift_mem_unitary}
    \lean{Matrix.weylClock_mem_unitary}
    \lean{Matrix.weyl_mem_unitary}
    \leanok
    Fix a dimension $d\geq 1$ and a primitive $d$-th root of unity $\zeta$.
    The cyclic shift $X$, the clock operator $Z$, and every Weyl operator
    $W(a,b)=X^aZ^b$ are unitary.
\end{theorem}

\begin{proof}\leanok
    The cyclic shift is the permutation matrix of a cyclic permutation.  The
    clock operator is diagonal, and each diagonal entry is a power of $\zeta$,
    hence has modulus one.  Thus $X$ and $Z$ are unitary, and so is every product
    $X^aZ^b$.
\end{proof}

\begin{theorem}[Weyl-operator unitary 1-design twirl]
    \label{thm:weyl_twirl}
    \lean{Matrix.sum_weyl_conj}
    \leanok
    \uses{lem:root_of_unity_orthogonality}
    Fix a dimension $d \ge 1$ and a primitive $d$-th root of unity $\zeta$, and
    let $X$ be the cyclic shift $\lvert i\rangle \mapsto \lvert i+1\rangle$ and
    $Z = \diag(\zeta^0,\dots,\zeta^{d-1})$ the clock operator. For
    every matrix $M$ on $\mathbb{C}^d$, the uniform average of the conjugations
    by the $d^2$ Weyl operators $W(a,b) = X^a Z^b$ is the completely depolarizing
    channel:
    \[
        \frac{1}{d^2} \sum_{a,b=0}^{d-1} W(a,b)\,M\,W(a,b)^{\dagger}
            = \frac{\tr M}{d}\,\mathbf 1.
    \]
\end{theorem}

\begin{proof}\leanok
    \uses{lem:root_of_unity_orthogonality}
    The double average factors into a clock average followed by a shift average.
    The clock average $\sum_{b} Z^{b} M (Z^{b})^{\dagger}$ multiplies the entry
    $M_{ij}$ by $\sum_{b} \zeta^{b\,i}\,\overline{\zeta^{b\,j}}$, which by
    Lemma~\ref{lem:root_of_unity_orthogonality} is $d$ when $i = j$ and $0$
    otherwise; the result is $d$ times the diagonal part of $M$. The shift
    average $\sum_{a} X^{a} (\operatorname{diagonal} v) (X^{a})^{\dagger}$
    cyclically permutes the diagonal entries, so each diagonal position receives
    the full sum $\sum_{k} v_k = \tr M$, giving
    $(\tr M)\,\mathbf 1$. Combining the two factors of $d$ with the
    prefactor $d^{-2}$ leaves $(\tr M / d)\,\mathbf 1$.
\end{proof}

\begin{lemma}[Twirl as partial trace tensored with the maximally mixed ancilla]
    \label{lem:twirl_partial_trace}
    \lean{Matrix.sum_kronecker_one_weyl_conj}
    \leanok
    \uses{thm:weyl_twirl}
    Fix a dimension $d_C \ge 1$ and a primitive $d_C$-th root of unity $\zeta$.
    For every matrix $M$ on $\mathcal H_S \otimes \mathbb{C}^{d_C}$, the uniform
    average of the conjugations by the $d_C^2$ unitaries
    $\mathbf 1_S \otimes W(a,b)$ on the second factor is the partial trace over
    that factor tensored with the maximally mixed state $\mathbf 1_C / d_C$:
    \[
        \frac{1}{d_C^2} \sum_{a,b}
            (\mathbf 1_S \otimes W(a,b))\, M\, (\mathbf 1_S \otimes W(a,b))^{\dagger}
            = (\tr_C M) \otimes (\mathbf 1_C / d_C).
    \]
\end{lemma}

\begin{proof}\leanok
    \uses{thm:weyl_twirl}
    On each pair of blocks indexed by the first factor, the conjugation by
    $\mathbf 1_S \otimes W(a,b)$ acts as the Weyl conjugation $W(a,b)\,(\cdot)\,
    W(a,b)^{\dagger}$ of the corresponding block of $M$. Averaging over the $d_C^2$
    Weyl operators sends each block to the depolarizing channel by
    Theorem~\ref{thm:weyl_twirl}, replacing it by its trace times
    $\mathbf 1_C / d_C$. The block trace is exactly the corresponding entry of the
    partial trace $\tr_C M$, so the average is $(\tr_C M) \otimes
    (\mathbf 1_C / d_C)$.
\end{proof}

\begin{theorem}[The maximally mixed ancilla]
    \label{thm:maximally_mixed_ancilla}
    \lean{Matrix.maximallyMixed_posDef}
    \lean{Matrix.maximallyMixed_trace}
    \leanok
    For $d_C\geq 1$, the matrix $\tau_C=d_C^{-1}\mathbf 1_C$ is positive
    definite and has trace one.
\end{theorem}

\begin{proof}\leanok
    The identity is positive definite and $d_C^{-1}>0$, so $\tau_C$ is positive
    definite.  Moreover,
    $\tr\tau_C=d_C^{-1}\tr\mathbf 1_C=1$.
\end{proof}

\begin{lemma}[Reindexing invariance of the quantum relative entropy]
    \label{lem:relative_entropy_submatrix_equiv}
    \lean{Matrix.quantumRelativeEntropy_submatrix_equiv}
    \leanok
    \uses{def:quantum_relative_entropy}
    For Hermitian matrices $\rho, \sigma$ on a finite index set and any bijection
    $e$ from that set onto another finite set,
    \[
        D(\rho_{e^{-1},e^{-1}} \,\Vert\, \sigma_{e^{-1},e^{-1}}) = D(\rho \,\Vert\, \sigma).
    \]
\end{lemma}

\begin{proof}\leanok
    \uses{def:quantum_relative_entropy}
    Write $D(\rho \,\Vert\, \sigma) = \tr\bigl(\rho(\log\rho - \log\sigma)\bigr)$.
    The matrix logarithm is covariant under reindexing,
    $\log(\rho_{e^{-1},e^{-1}}) = (\log\rho)_{e^{-1},e^{-1}}$, because reindexing is a
    star-algebra isomorphism and so commutes with the continuous functional calculus.
    The same isomorphism preserves products and the trace,
    \[
        (\rho\,M)_{e^{-1},e^{-1}} = \rho_{e^{-1},e^{-1}}\,M_{e^{-1},e^{-1}},
        \qquad
        \tr\bigl(M_{e^{-1},e^{-1}}\bigr) = \tr(M).
    \]
    Taking $M = \log\rho - \log\sigma$ and applying these three identities gives
    $D(\rho_{e^{-1},e^{-1}} \,\Vert\, \sigma_{e^{-1},e^{-1}}) = D(\rho \,\Vert\, \sigma)$.
\end{proof}

\begin{theorem}[Data-processing inequality under the partial trace]
    \label{thm:relative_entropy_data_processing}
    \lean{Matrix.quantumRelativeEntropy_partialTraceRight_le}
    \leanok
    \uses{def:quantum_relative_entropy, thm:relative_entropy_ancilla_additivity,
        thm:relative_entropy_unitary_invariance, thm:relative_entropy_joint_convexity,
        lem:relative_entropy_submatrix_equiv, lem:twirl_partial_trace}
    For positive definite matrices $\rho, \sigma$ on a tensor product of a system
    factor and an ancilla factor,
    \[
        D(\tr_C \rho \,\Vert\, \tr_C \sigma) \le D(\rho \,\Vert\, \sigma),
    \]
    where $\tr_C$ is the partial trace over the ancilla factor. This is the
    positive-definite base case; the source inequality on the support domain
    $\ker \sigma \subseteq \ker \rho$ is
    Theorem~\ref{thm:relative_entropy_data_processing_support}.
\end{theorem}

\begin{proof}\leanok
    \uses{thm:relative_entropy_ancilla_additivity,
        thm:relative_entropy_unitary_invariance, thm:relative_entropy_joint_convexity,
        lem:twirl_partial_trace}
    By ancilla additivity (Theorem~\ref{thm:relative_entropy_ancilla_additivity})
    the reduced-state relative entropy equals
    $D\bigl((\tr_C \rho) \otimes (\mathbf 1_C / d_C) \,\Vert\,
    (\tr_C \sigma) \otimes (\mathbf 1_C / d_C)\bigr)$. By
    Lemma~\ref{lem:twirl_partial_trace} each tensored reduced state is the uniform
    average of the conjugations by the $d_C^2$ unitaries
    $U_{ab} = \mathbf 1_S \otimes W(a,b)$, so this is the relative entropy of a
    convex combination of the conjugated pairs
    $(U_{ab}\,\rho\,U_{ab}^{\dagger}, U_{ab}\,\sigma\,U_{ab}^{\dagger})$ with equal
    weights $d_C^{-2}$. Joint convexity bounds it above by the same convex
    combination of the per-term relative entropies, each of which equals
    $D(\rho \,\Vert\, \sigma)$ by unitary invariance
    (Theorem~\ref{thm:relative_entropy_unitary_invariance}). As the weights sum to
    one, the bound is $D(\rho \,\Vert\, \sigma)$.
\end{proof}

\begin{theorem}[Data-processing inequality on the support domain]
    \label{thm:relative_entropy_data_processing_support}
    \lean{Matrix.quantumRelativeEntropy_partialTraceRight_le_support}
    \leanok
    \uses{def:quantum_relative_entropy, thm:relative_entropy_data_processing}
    For positive semidefinite $\rho, \sigma$ on a tensor product of a system factor
    and an ancilla factor of dimension $d_C$, with the support condition
    $\ker \sigma \subseteq \ker \rho$,
    \[
        D(\tr_C \rho \,\Vert\, \tr_C \sigma) \le D(\rho \,\Vert\, \sigma),
    \]
    where $\tr_C$ is the partial trace over the ancilla factor.
\end{theorem}

\begin{proof}\leanok
    \uses{thm:relative_entropy_data_processing}
    Let $N$ be the dimension of the joint space. Regularize both arguments through
    the affine trace-shrinking perturbation
    $M_\varepsilon = (1+N\varepsilon)^{-1}(M+\varepsilon\,\mathbf 1)$, which is
    positive definite for every $\varepsilon>0$, so the positive-definite
    data-processing inequality
    (Theorem~\ref{thm:relative_entropy_data_processing}) gives
    \[
        D(\tr_C \rho_\varepsilon \,\Vert\, \tr_C \sigma_\varepsilon)
            \le D(\rho_\varepsilon \,\Vert\, \sigma_\varepsilon),
        \qquad \varepsilon > 0.
    \]
    The right-hand side converges to $D(\rho \,\Vert\, \sigma)$ as
    $\varepsilon\to0^+$, by the same shared-eigenbasis scalar-limit argument as the
    joint convexity on the support domain
    (Theorem~\ref{thm:relative_entropy_joint_convexity_support}).

    For the left-hand side, the partial trace of the regularization is a
    differently-scaled regularization of the marginal: because the partial trace is
    linear and $\tr_C \mathbf 1 = d_C\,\mathbf 1$,
    \[
        \tr_C\bigl((1+N\varepsilon)^{-1}(M+\varepsilon\,\mathbf 1)\bigr)
            = (1+N\varepsilon)^{-1}\bigl(\tr_C M + d_C\varepsilon\,\mathbf 1\bigr),
    \]
    the affine regularization of $\tr_C M$ with the same scaling rate $N$ but shift
    rate $d_C$ and the smaller identity on the system factor. The support condition
    transfers to the marginals,
    $\ker(\tr_C \sigma) \subseteq \ker(\tr_C \rho)$: a vector annihilated by
    $\tr_C \sigma$ has vanishing marginal quadratic form, which splits into the
    nonnegative joint quadratic forms of its single-ancilla lifts, so each lift lies
    in $\ker \sigma$, hence in $\ker \rho$, and reassembling the ancilla sum shows
    the vector lies in $\ker(\tr_C \rho)$. With this support condition the
    arbitrary-rate affine regularization has the same both-arguments limit, so
    \[
        D(\tr_C \rho_\varepsilon \,\Vert\, \tr_C \sigma_\varepsilon)
            \to D(\tr_C \rho \,\Vert\, \tr_C \sigma),
        \qquad \varepsilon\to0^+.
    \]
    Passing the inequality through the two limits gives the support-domain bound.
\end{proof}
